\section{Data Preprocessing}

The AI-Assisted Maternal Health Risk Screening System is trained using one
dataset, \texttt{Mathernal\_Risk.csv}. It contains 6,103 records and three
target classes: Low risk, Mid risk, and High risk. The model uses eight
numerical features: age, body temperature, heart rate, systolic blood
pressure, diastolic blood pressure, body mass index (BMI), HbA1c, and
fasting glucose. Patient names and identifiers are excluded from the model
features.

\subsection{Data Validation and Cleaning}

Before training, the dataset is checked for the expected columns, valid
risk labels, and numeric feature values. Exact duplicate records are
removed. Records outside configured physiological plausibility ranges are
also removed; these ranges are intended to identify implausible values,
not to act as clinical thresholds.

Near-duplicate records are grouped using their feature values rounded to
whole numbers. Records in the same group are kept together when the data
is divided into training, validation, and test sets. This reduces the
chance that near-identical records appear in both training and test data
and inflate the measured performance.

\subsection{Missing Values and Feature Scaling}

The supplied dataset contains no missing feature values. The preprocessing
pipeline includes median imputation, with medians fitted using the training
data only and saved for use during prediction. Imputation is not performed
separately for each risk class.

Min--Max normalization is not applied to all features. Tree-based models
use the measurements without scaling. Logistic Regression uses
standardization within its model pipeline. The same saved preprocessing
steps are applied to new inputs.

\subsection{Class Labels and Class Imbalance}

The target labels are encoded consistently as Low risk, Mid risk, and
High risk. The supplied dataset is approximately balanced across the
three classes. SMOTE is not used. Class weights are used by selected
classifiers.

No additional derived clinical features are introduced for model training.
The system uses the eight measurements listed above as model inputs.
Symptoms, assessment notes, and patient-record information are handled
separately by the system's rules and advisory workflows.

\subsection{Data Limitations}

The dataset's risk labels are supplied by the source dataset, but the
criteria and process used to assign those labels are undocumented. In
addition, the glucose column headers do not consistently match the
magnitudes of their values. The values are used as supplied, and the
dataset validation checks their expected ranges. These limitations should
be considered when interpreting model results.

\section{Evaluation Metrics}

The training pipeline compares six classifiers: Random Forest, XGBoost,
Decision Tree, Extra Trees, Gradient Boosting, and Logistic Regression.
The cleaned data is divided into approximately 70\% training, 15\%
validation, and 15\% test data using a group-aware, approximately
stratified split. The validation data is used to select the production
model. The held-out test data is used to report model performance.

\subsection{Classification Performance Metrics}

Model performance is evaluated using accuracy, macro-averaged precision,
macro-averaged recall, macro-averaged F1-score, High-risk recall, a
confusion matrix, and one-vs-rest macro ROC-AUC. Macro-averaging gives
equal weight to each risk class. High-risk recall is reported separately
to show how many records labelled High risk in the test data were
identified as High risk by the model.

Table~\ref{tab:model-performance} shows the test-set results from the
current training run. XGBoost was selected as the production model using
validation metrics, with High-risk recall prioritized in model selection.

\begin{table}[htbp]
\centering
\small
\caption{Classifier performance on the held-out test split}
\label{tab:model-performance}
\begin{tabular}{lrrrrr}
\hline
Model & Accuracy & Macro precision & Macro recall & Macro F1 & High-risk recall \\
\hline
Random Forest       & 99.01\% & 99.02\% & 99.01\% & 99.01\% & 98.03\% \\
XGBoost             & 99.78\% & 99.78\% & 99.78\% & 99.78\% & 99.67\% \\
Decision Tree       & 97.80\% & 97.82\% & 97.81\% & 97.81\% & 96.05\% \\
Extra Trees         & 97.36\% & 97.40\% & 97.37\% & 97.36\% & 94.41\% \\
Gradient Boosting   & 99.23\% & 99.23\% & 99.23\% & 99.23\% & 98.36\% \\
Logistic Regression & 57.32\% & 56.98\% & 57.43\% & 56.41\% & 45.39\% \\
\hline
\end{tabular}
\end{table}

\subsection{Prediction Probabilities and Explanations}

The model returns a probability for each risk class. The system's
confidence score is the largest of the three class probabilities:

\begin{equation}
    \text{Confidence} =
    \max(P_{\text{Low}},\, P_{\text{Mid}},\, P_{\text{High}})
    \label{eq:confidence}
\end{equation}

This score represents the model's predicted probability and should not
be interpreted as medical certainty or a diagnosis. The training
evaluation also calculates probability-calibration measures, including
the multiclass Brier score and expected calibration error.

The system provides feature-contribution explanations. Tree-based models
use SHAP where available, while Logistic Regression uses linear-model
contributions. These explanations describe how features contributed to
the model's prediction; they do not establish medical causes.

\subsection{Rules, Alerts, and Advisory Generation}

The system's clinical rules engine evaluates assessment measurements and
selected symptoms separately from the machine-learning classifier. Its
alerts are distinct from the model's predicted risk class.

The advisory workflow uses relevant assessment and patient-record
information together with passages retrieved from the system's knowledge
base. It uses a locally configured Ollama model when available and has an
offline knowledge-based fallback. The system does not use Gemini to
generate advisories.

\subsection{Interpretation and Limitations of Results}

The metrics in Table~\ref{tab:model-performance} measure agreement with
the supplied dataset's labels on an internal held-out test split. Because
the source and clinical criteria for those labels are undocumented, these
results do not establish clinical validity, performance on other
populations, or suitability for clinical decision-making. External
evaluation using clinically reviewed data would be needed to assess
those properties.

Cosine similarity for advisory text, a combined System Performance Score
(SPS), and a Referral Match Score (RMS) are not currently calculated by
the system. They should only be described as future evaluation measures
if they are planned for later implementation.
